\ifdefined\gkver\else\def\gkver{student}\fi
\documentclass[\gkver]{gaokaozhenti}
\begin{document}

\chapter{2017年上海}

\begin{enumerate}
\item
%% number: 1
%% paperName: 2017年上海等级考第 1 题 3 分
%% typeId: 1
%% type: 单选题
%% chapter: 原子和原子核
%% point: 天然放射性
%% method: 无
%% score: 3
%% degree: 950
%% duplicateId: 0
%% body:
由放射性元素放出的氦核流被称为\xzanswer{B}
\fourchoices[answer=B]
{阴极射线}
{$\alpha$ 射线}
{$\beta$ 射线}
{$\gamma$ 射线}
%% answer: B
%% memo:
\memoanswer{
本题原卷及材料中未提供详解，待补充。
}

\item
%% number: 2
%% paperName: 2017年上海等级考第 2 题 3 分
%% typeId: 1
%% type: 单选题
%% chapter: 光学
%% point: 光子说
%% method: 无
%% score: 3
%% degree: 950
%% duplicateId: 0
%% body:
光子的能量与其\xzanswer{A}
\fourchoices[answer=A]
{频率成正比}
{波长成正比}
{速度成正比}
{速度平方成正比}
%% answer: A
%% memo:
\memoanswer{
本题原卷及材料中未提供详解，待补充。
}

\item
%% number: 3
%% paperName: 2017年上海等级考第 3 题 3 分
%% typeId: 1
%% type: 单选题
%% chapter: 原子和原子核
%% point: 原子核的组成
%% method: 无
%% score: 3
%% degree: 950
%% duplicateId: 0
%% body:
在同位素氢、氘，氚的核内具有相同的\xzanswer{D}
\fourchoices[answer=D]
{核子数}
{电子数}
{中子数}
{质子数}
%% answer: D
%% memo:
\memoanswer{
本题原卷及材料中未提供详解，待补充。
}

\item
%% number: 4
%% paperName: 2017年上海等级考第 4 题 3 分
%% typeId: 1
%% type: 单选题
%% chapter: 光学
%% point: 光的干涉
%% method: 无
%% score: 3
%% degree: 850
%% duplicateId: 0
%% body:
用单色光照射位于竖直平面内的肥皂液薄膜，所观察到的干涉条纹为\xzanswer{B}
\fourchoices[answer=B, ispicture=true, h=2.6cm]
{figs/04a.png}
{figs/04b.png}
{figs/04c.png}
{figs/04d.png}
%% answer: B
%% memo:
\memoanswer{
本题原卷及材料中未提供详解，待补充。
}

\item
%% number: 5
%% paperName: 2017年上海等级考第 5 题 3 分
%% typeId: 1
%% type: 单选题
%% chapter: 电场
%% point: 带电粒子的直线运动
%% method: 无
%% score: 3
%% degree: 850
%% duplicateId: 0
%% body:
如图，在匀强电场中，悬线一端固定于地面，另一端拉住一个带电小球，使之处于静止状态。忽略空气阻力，当悬线断裂后，小球将做\xzanswer{C}
\onepicture[h=3.2cm, align=right]{\includesvg{figs/05}}
\fourchoices[answer=C]
{曲线运动}
{匀速直线运动}
{匀加速直线运动}
{变加速直线运动}
%% answer: C
%% memo:
\memoanswer{
本题原卷及材料中未提供详解，待补充。
}

\item
%% number: 6
%% paperName: 2017年上海等级考第 6 题 3 分
%% typeId: 1
%% type: 单选题
%% chapter: 运动和力
%% point: 牛顿第二定律
%% method: 无
%% score: 3
%% degree: 850
%% duplicateId: 0
%% body:
一碗水置于火车车厢内的水平桌面上。当火车向右做匀减速运动时，水面形状接近于图\xzanswer{A}
\fourchoices[answer=A, ispicture=true, h=2.2cm]
{figs/06a.png}
{figs/06b.png}
{figs/06c.png}
{figs/06d.png}
%% answer: A
%% memo:
\memoanswer{
本题原卷及材料中未提供详解，待补充。
}

\item
%% number: 7
%% paperName: 2017年上海等级考第 7 题 3 分
%% typeId: 1
%% type: 单选题
%% chapter: 功和能
%% point: 动能势能
%% method: 近似与估算
%% score: 3
%% degree: 850
%% duplicateId: 0
%% body:
从大型加速器射出的电子束总能量约为 $1\UGeV$（$1\UGeV = 1.6\times10^{-10}\UJ$），此能量最接近\xzanswer{A}
\fourchoices[answer=A]
{一只爬行的蜗牛的动能}
{一个奔跑的孩子的动能}
{一辆行驶的轿车的动能}
{一架飞行的客机的动能}
%% answer: A
%% memo:
\memoanswer{
本题原卷及材料中未提供详解，待补充。
}

\item
%% number: 8
%% paperName: 2017年上海等级考第 8 题 3 分
%% typeId: 1
%% type: 单选题
%% chapter: 气体性质
%% point: 物体的内能
%% method: 无
%% score: 3
%% degree: 750
%% duplicateId: 0
%% body:
一个密闭容器由固定导热板分隔为体积相同的两部分，分别装有质量不等的同种气体。当两部分气体稳定后，它们的\xzanswer{C}
\fourchoices[answer=C]
{密度相同}
{分子数相同}
{分子平均速率相同}
{分子间平均距离相同}
%% answer: C
%% memo:
\memoanswer{
由于是导热板，因此两部分气体的温度最终会相同，而温度是分子平均动能的标志，所以分子平均动能也相同。

高中课本只讲平均动能（大学物理的结论为 $\frac{3}{2}kT$），不讲平均速率。平均动能相等时平均速率一定相等吗？按照高中生现有的知识，应该这样做：在同种气体的情况下，$m$ 相同，即证明\textbf{在 $v^{2}$ 的平均值相等的情况下，$v$ 的平均值也相等}，显然\textbf{这个结论在数学上是错误的}！真正的证明需要大学物理，可证平均速率为 $\sqrt{\frac{8kT}{\pi m}}$，这样才能得出平均速率也不变的结论。
}

\item
%% number: 9
%% paperName: 2017年上海等级考第 9 题 4 分
%% typeId: 1
%% type: 单选题
%% chapter: 电路
%% point: 电路中的图象
%% method: 无
%% score: 4
%% degree: 850
%% duplicateId: 0
%% body:
将四个定值电阻 a、b、c、d 分别接入电路，测得相应的电流、电压值如图所示。其中阻值最接近的两个电阻是\xzanswer{A}
\onepicture[h=3.2cm, align=right]{\includesvg{figs/09}}
\fourchoices[answer=A]
{a 和 b}
{b 和 d}
{a 和 c}
{c 和 d}
%% answer: A
%% memo:
\memoanswer{
本题原卷及材料中未提供详解，待补充。
}

\item
%% number: 10
%% paperName: 2017年上海等级考第 10 题 4 分
%% typeId: 1
%% type: 单选题
%% chapter: 机械振动
%% point: 单摆
%% method: 无
%% score: 4
%% degree: 650
%% duplicateId: 0
%% body:
做简谐运动的单摆，其摆长不变，若摆球的质量增加为原来的 $\frac{9}{4}$ 倍，摆球经过平衡位置的速度减为原来的 $\frac{2}{3}$，则单摆振动的\xzanswer{B}
\fourchoices[answer=B]
{周期不变，振幅不变}
{周期不变，振幅变小}
{周期改变，振幅不变}
{周期改变，振幅变大}
%% answer: B
%% memo:
\memoanswer{
本题原卷及材料中未提供详解，待补充。
}

\item
%% number: 11
%% paperName: 2017年上海等级考第 11 题 4 分
%% typeId: 1
%% type: 单选题
%% chapter: 磁场
%% point: 左手定则
%% method: 无
%% score: 4
%% degree: 650
%% duplicateId: 0
%% body:
如图，一导体棒 ab 静止在 U 型铁芯的两臂之间。电键闭合后导体棒受到的安培力方向\xzanswer{D}
\onepicture[h=3.4cm, align=right]{\includesvg{figs/11}}
\fourchoices[answer=D]
{向上}
{向下}
{向左}
{向右}
%% answer: D
%% memo:
\memoanswer{
本题原卷及材料中未提供详解，待补充。
}

\item
%% number: 12
%% paperName: 2017年上海等级考第 12 题 4 分
%% typeId: 1
%% type: 单选题
%% chapter: 气体性质
%% point: 液柱移动定性分析
%% method: 无
%% score: 4
%% degree: 550
%% duplicateId: 0
%% body:
如图，竖直放置的 U 形管内装有水银，左端开口，右端封闭一定量的气体，底部有一阀门。开始时阀门关闭，左管的水银面较高。现打开阀门，流出一些水银后关闭阀门。当重新平衡时\xzanswer{D}
\onepicture[h=4.2cm, align=right]{\includesvg{figs/12}}
\fourchoices[answer=D]
{左管的水银面与右管等高}
{左管的水银面比右管的高}
{左管的水银面比右管的低}
{水银面高度关系无法判断}
%% answer: D
%% memo:
\memoanswer{
本题原卷及材料中未提供详解，待补充。
}

\item
%% number: 13
%% paperName: 2017年上海等级考第 13 题 4 分
%% typeId: 3
%% type: 填空题
%% chapter: 电场
%% point: 电势能电势差电场力做功
%% method: 无
%% score: 4
%% degree: 850
%% duplicateId: 0
%% body:
静电场中某电场线如图所示。把点电荷从电场中的 A 点移到 B 点，其电势能增加 $1.2\times10^{-7}\UJ$，则该点电荷带\tkanswer[1.2cm]{负}电（选填：“正”或“负”）；在此过程中电场力做功为\tkanswer[2.6cm]{$-1.2\times10^{-7}$}$\UJ$。
\onepicture[h=2.4cm, align=right]{\includesvg{figs/13}}
%% answer: 负，$-1.2\times10^{-7}$
%% memo:
\memoanswer{
本题原卷及材料中未提供详解，待补充。
}

\item
%% number: 14
%% paperName: 2017年上海等级考第 14 题 4 分
%% typeId: 3
%% type: 填空题
%% chapter: 机械波
%% point: 波长频率波速
%% method: 无
%% score: 4
%% degree: 850
%% duplicateId: 0
%% body:
机械波产生和传播的条件是：①存在一个做振动的波源，②在波源周围存在\tkanswer[2.2cm]{介质}；机械波传播的是\tkanswer[2.6cm]{运动形式}和\tkanswer[2.6cm]{能量（或信息）}。
%% answer: 介质，运动形式，能量（或信息）
%% memo:
\memoanswer{
本题原卷及材料中未提供详解，待补充。
}

\item
%% number: 15
%% paperName: 2017年上海等级考第 15 题 4 分
%% typeId: 3
%% type: 填空题
%% chapter: 直线运动
%% point: 竖直上下抛运动
%% method: 无
%% score: 4
%% degree: 750
%% duplicateId: 0
%% body:
物体以 $25\Ums$ 的初速度做竖直上抛运动，经过\tkanswer[1.6cm]{2.5}$\Us$ 到达最高点，它在第三秒内的位移为\tkanswer[1.6cm]{0}$\Um$。（$g$ 取 $10\Umsq$）
%% answer: 2.5，0
%% memo:
\memoanswer{
本题原卷及材料中未提供详解，待补充。
}

\item
%% number: 16
%% paperName: 2017年上海等级考第 16 题 4 分
%% typeId: 3
%% type: 填空题
%% chapter: 气体性质
%% point: 查理定律
%% method: 无
%% score: 4
%% degree: 750
%% duplicateId: 0
%% body:
如图，气缸固定于水平面，用截面积为 $20\Ucmq$ 的活塞封闭一定量的气体，活塞与缸壁间摩擦不计。当大气压强为 $1.0\times10^{5}\UPa$、气体温度为 $87\UdC$ 时，活塞在大小为 $40\UN$、方向向左的力 $F$ 作用下保持静止，气体压强为\tkanswer[3.0cm]{$1.2\times10^{5}$}$\UPa$。若保持活塞不动，将气体温度降至 $27\UdC$，则 $F$ 变为\tkanswer[1.6cm]{0}$\UN$。
\onepicture[h=2.4cm, align=right]{\includesvg{figs/16}}
%% answer: $1.2\times10^{5}$，0
%% memo:
\memoanswer{
（1）对活塞进行受力分析，有：$p_{1}S = F + p_{0}S$，解得 $p_{1} = p_{0} + \frac{F}{S} = 1.0\times10^{5} + \frac{40}{20\times10^{-4}} = 1.2\times10^{5}\UPa$；

（2）活塞不动，气体体积不变，由查理定律：$\frac{p_{1}}{T_{1}} = \frac{p_{2}}{T_{2}}$，$\frac{1.2\times10^{5}}{360} = \frac{p_{2}}{300}$，解得 $p_{2} = 1.0\times10^{5}\UPa$。

对活塞进行受力分析，有：$p_{2}S = F' + p_{0}S$，解得 $F' = 0$。
}

\item
%% number: 17
%% paperName: 2017年上海等级考第 17 题 4 分
%% typeId: 3
%% type: 填空题
%% chapter: 功和能
%% point: 机械能守恒定律
%% method: 无
%% score: 4
%% degree: 650
%% duplicateId: 0
%% body:
如图，光滑固定斜面的倾角为 $30\Udu$，A、B 两物体的质量之比为 $4:1$。B 用不可伸长的轻绳分别与 A 和地面相连，开始时 A、B 离地高度相同。在 C 处剪断轻绳，当 B 落地前瞬间，A、B 的速度大小之比为\tkanswer[1.6cm]{$1:2$}，机械能之比为\tkanswer[1.6cm]{$4:1$}（以地面为零势能面）。
\onepicture[h=2.6cm, align=right]{\includesvg{figs/17}}
%% answer: $1:2$，$4:1$
%% memo:
\memoanswer{
本题原卷及材料中未提供详解，待补充。
}

\item
%% number: 18
%% paperName: 2017年上海等级考第 18 题 10 分
%% typeId: 4
%% type: 实验题
%% chapter: 电路
%% point: 测电动势与内阻
%% method: 无
%% score: 10
%% degree: 650
%% duplicateId: 0
%% body:
“用 DIS 测定电源的电动势和内阻”的实验电路如图~\ref{2017上海18a} 所示，其中定值电阻阻值 $R_{1} = 1\UO$。
\begin{enumerate}
	\item
	图~\ref{2017上海18a} 中 A 为\tkanswer[1.8cm]{电流}传感器，定值电阻 $R_{1}$ 在实验中起\tkanswer[2.8cm]{保护电路}的作用；
	\item
	实验测得的路端电压 $U$ 相应电流 $I$ 的拟合曲线如图~\ref{2017上海18b} 所示，由此得到电源电动势 $E =$\tkanswer[1.6cm]{1.50}$\UV$，内阻 $r =$\tkanswer[1.6cm]{0.28}$\UO$；
	\item
	实验测得的数据如表所示，则实验中选用的滑动变阻器最合理的阻值范围为（\quad）

	\begin{center}
	\begin{tabular}{|c|c|}
	\hline
	$U/\UV$ & $I/\UA$ \\ \hline
	1.48 & 0.106 \\ \hline
	1.44 & 0.202 \\ \hline
	1.39 & 0.411 \\ \hline
	1.34 & 0.561 \\ \hline
	1.23 & 0.948 \\ \hline
	\end{tabular}
	\end{center}

	\fourchoices[answer=B]
	{$0\sim5\UO$}
	{$0\sim20\UO$}
	{$0\sim50\UO$}
	{$0\sim200\UO$}
\end{enumerate}
\twopicture[labelA=2017上海18a, labelB=2017上海18b, numA=（a）, numB=（b）, hA=4.2cm, hB=4.2cm, align=right]
{figs/18a.png}
{figs/18b.png}
%% answer: （1）电流，保护电路（2）1.50，0.28（3）B
\jdanswer{
\begin{enumerate}
	\item
	电流，保护电路

	\item
	1.50，0.28

	\item
	B

\end{enumerate}
}
%% memo:
\memoanswer{
由表格第 1 行数据可知，此实验外电阻最大调节到了 $\frac{1.48}{0.106} \approx 14\UO$，因此应选用 $0\sim20\UO$ 的滑动变阻器，而滑动变阻器阻值太大会使实验调节不便，所以 C、D 选项错误。正确选项为 B。
}

\item
%% number: 19
%% paperName: 2017年上海等级考第 19 题 14 分
%% typeId: 6
%% type: 计算题
%% chapter: 运动和力
%% point: 牛顿定律的应用
%% method: 无
%% score: 14
%% degree: 550
%% duplicateId: 0
%% body:
如图，与水平面夹角 $\theta = 37\Udu$ 的斜面和半径 $R = 0.4\Um$ 的光滑圆轨道相切于 B 点，且固定于竖直平面内。滑块从斜面上的 A 点由静止释放，经 B 点后沿圆轨道运动，通过最高点 C 时轨道对滑块的弹力为零。已知滑块与斜面间动摩擦因数 $\mu = 0.25$。（$g$ 取 $10\Umsq$，$\sin37\Udu = 0.6$，$\cos37\Udu = 0.8$）求：
\begin{enumerate}
	\item
	滑块在 C 点的速度大小 $v_{C}$；
	\item
	滑块在 B 点的速度大小 $v_{B}$；
	\item
	A、B 两点间的高度差 $h$。
\end{enumerate}
\onepicture[w=7.5cm, align=right]{\includesvg{figs/19}}
%% answer: 
\jdanswer{
\begin{enumerate}
	\item
	$v_{C} = 2\Ums$

	\item
	$v_{B} = 4.29\Ums$

	\item
	$h = 1.38\Um$

\end{enumerate}
}
%% memo:
\memoanswer{
（1）由题意，在 C 处滑块仅在重力作用下做圆周运动。设滑块质量为 $m$，由牛顿第二定律

$mg = m\frac{v_{C}^{2}}{R}$

得

$v_{C} = \sqrt{gR} = 2\Ums$

（2）由几何关系，B、C 间高度差 $H$ 为

$H = R(1 + \cos37\Udu) = 0.72\Um$

滑块由 B 到 C 的运动过程中仅重力做功，机械能守恒。以 B 为势能零点

$\frac{1}{2}mv_{B}^{2} = mgH + \frac{1}{2}mv_{C}^{2}$

$v_{B} = \sqrt{2gH + v_{C}^{2}}$

代入已知数据得

$v_{B} = 4.29\Ums$

（3）滑块由 A 到 B 过程中受力如图

\onepicture[h=2.4cm]{\includesvg{figs/19_an}}

由牛顿定律

$mg\sin37\Udu - f = ma$

$N = mg\cos37\Udu$

$f = \mu N$

解得

$a = g(\sin37\Udu - \mu\cos37\Udu)$

代入已知数据得

$a = 4\Umsq$

设 A、B 间距为 $L$，由匀加速运动公式 $v_{B}^{2} = 2aL$

及几何关系

$h = L\sin37\Udu$

得

$h = \frac{v_{B}^{2}}{2a}\sin37\Udu = 1.38\Um$

19 题评分

\begin{center}
\begin{tabular}{|c|p{10cm}|}
\hline
评分等级 & 描述 \\ \hline
A（14分） & 从基本原理出发，说明清晰，推理合乎逻辑，演算准确。 \\ \hline
B（10-13分） & 有些未从基本原理出发，或说明推理不够理想，或部分演算有误。……。 \\ \hline
C（5-9分） & 约一半左右没有做出。或尽管做出大部分但解题思路混乱，…… \\ \hline
D（1-4分） & 做出一小部分，或只写出零星的一些公式与答案，…… \\ \hline
E（0分） & 没有一点合理，或没做。 \\ \hline
\end{tabular}
\end{center}
}

\item
%% number: 20
%% paperName: 2017年上海等级考第 20 题 16 分
%% typeId: 6
%% type: 计算题
%% chapter: 电磁感应
%% point: 电磁感应中的变加速
%% method: 无
%% score: 16
%% degree: 450
%% duplicateId: 0
%% body:
如图，光滑平行金属导轨间距为 $L$，与水平面夹角为 $\theta$，两导轨上端用阻值为 $R$ 的电阻相连，该装置处于磁感应强度为 $B$ 的匀强磁场中，磁场方向垂直于导轨平面。质量为 $m$ 的金属杆 ab 以沿导轨平面向上的初速度 $v_{0}$ 从导轨底端开始运动，然后又返回到出发位置。在运动过程中，ab 与导轨垂直且接触良好，不计 ab 和导轨的电阻及空气阻力。
\begin{enumerate}
	\item
	求 ab 开始运动时的加速度 $a$；
	\item
	分析并说明 ab 在整个运动过程中速度、加速度的变化情况；
	\item
	分析并比较 ab 上滑时间和下滑时间的长短。
\end{enumerate}
\onepicture[w=8.5cm, align=right]{\includesvg{figs/20}}
%% answer: 
\jdanswer{
\begin{enumerate}
	\item
	$a = g\sin\theta + \frac{B^{2}L^{2}v_{0}}{mR}$

	\item
	杆上滑时：杆做加速度减小的减速运动，到达一定高度后速度为零。在最高点：杆的速度为零，加速度为 $g\sin\theta$，方向沿斜面向下。杆下滑时：杆做初速为零、加速度减小的加速运动。

	\item
	上滑所需时间小于下滑所需时间。

\end{enumerate}
}
%% memo:
\memoanswer{
（1）ab 开始运动时产生的感应电动势

$E = BLv_{0}$

回路中的感应电流

$I = \frac{E}{R}$

杆所受安培力

$F_{\mathrm{A}} = BIL = \frac{B^{2}L^{2}v_{0}}{R}$

杆受力如图，由牛顿定律

\onepicture[h=2.6cm]{\includesvg{figs/20_an}}

$mg\sin\theta + F_{\mathrm{A}} = ma$

得

$a = g\sin\theta + \frac{B^{2}L^{2}v_{0}}{mR}$

（2）分析如下：

杆上滑时：合力 $F = mg\sin\theta + F_{\mathrm{A}}$，与运动方向相反，杆减速；随着速度减小，$F_{\mathrm{A}}$ 减小，合力减小，加速度减小；因此杆做加速度减小的减速运动，到达一定高度后速度为零。

在最高点：杆的速度为零，加速度为 $g\sin\theta$，方向沿斜面向下。

杆下滑时：合力 $F = mg\sin\theta - F_{\mathrm{A}}$，与运动方向相同，杆加速；随着速度增加，$F_{\mathrm{A}}$ 增大，合力减小，加速度减小；因此杆做初速为零、加速度减小的加速运动。

（3）分析如下：

上滑和下滑过程经过的位移大小相等。而上滑时杆加速度大于 $g\sin\theta$，下滑时杆加速度小于 $g\sin\theta$，因此上升时的平均加速度大于下降时的平均加速度，由运动学规律可知，上滑所需时间小于下滑所需时间。

20 题评分

\begin{center}
\begin{tabular}{|c|p{10cm}|}
\hline
评分等级 & 描述 \\ \hline
A（16分） & 从基本原理出发，说明清晰，推理合乎逻辑，演算准确。 \\ \hline
B（11-15分） & 有些未从基本原理出发，或说明推理不够理想，或部分演算有误。……。 \\ \hline
C（6-10分） & 约一半左右没有做出。或尽管做出大部分但解题思路混乱，…… \\ \hline
D（1-5分） & 做出一小部分，或只写出零星的一些公式与答案，…… \\ \hline
E（0分） & 没有一点合理，或没做。 \\ \hline
\end{tabular}
\end{center}
}

\end{enumerate}

\end{document}
